% Formalization sections for P-A-R: A Type-Safe Agent Runtime in OCaml
% ICFP Workshop / ML Workshop submission

\documentclass[11pt]{article}

\usepackage{amsmath}
\usepackage{amssymb}
\usepackage{amsthm}
\usepackage{mathpartir}
\usepackage{listings}
\usepackage{xspace}

% Theorem styles
\theoremstyle{plain}
\newtheorem{prop}[equation]{Property}

% OCaml code listings
\lstdefinestyle{ocaml}{
  language=Caml,
  basicstyle=\ttfamily,
  keywordstyle=\bfseries,
  commentstyle=\itshape,
  morestring=[b]",
  stringstyle=\ttfamily,
  showstringspaces=false,
  breaklines=true,
  frame=single,
  numbers=left,
  numberstyle=\tiny,
  stepnumber=1,
  numbersep=5pt,
  tabsize=2,
  morecomment=[l]{(*},
  morecomment=[n]{(*}{*)},
}

% Math commands
\providecommand{\N}{\mathbb{N}}
\providecommand{\powerset}{\mathcal{P}}
\providecommand{\Transitions}{\longrightarrow}
\providecommand{\given}{\mid}
\DeclareRobustCommand{\identifiable}[1]{\texttt{#1}}
\providecommand{\iiff}{\Leftrightarrow}

%---------------------------------------------------------------------

\section{Formal Model}

%---------------------------------------------------------------------

\subsection{Task Lifecycle as a Labelled Transition System}
\label{sec:lts}

We model the task lifecycle as a labelled transition system (LTS) that captures
all valid states and transitions in the P-A-R runtime.

\subsubsection{Syntax}

The LTS is a tuple $(\mathcal{S}, \mathcal{L}, \rightarrow)$ where:

\begin{itemize}
  \item $\mathcal{S}$ is the set of states
  \item $\mathcal{L}$ is the set of transition labels
  \item $\rightarrow \subseteq \mathcal{S} \times \mathcal{L} \times \mathcal{S}$ is the transition relation
\end{itemize}

\begin{definition}[State Space]
  \label{def:states}
  \[
  \mathcal{S} \triangleq \{\, \identifiable{Pending},\ \identifiable{Scheduled},\ \identifiable{Running},\ \identifiable{Waiting\_input},\ \identifiable{Suspended},\ \identifiable{Completed},\ \identifiable{Failed},\ \identifiable{Cancelled} \,\}
  \]
\end{definition}

\begin{definition}[Terminal States]
  \label{def:terminal}
  \[
  \mathcal{T} \triangleq \{\, \identifiable{Completed},\ \identifiable{Failed},\ \identifiable{Cancelled} \,\}
  \]
  No outgoing transitions are defined from any $s \in \mathcal{T}$.
\end{definition}

\begin{definition}[Transition Labels]
  \label{def:labels}
  The set of labels $\mathcal{L}$ classifies the reasons for transitions:
  \[
  \mathcal{L} \triangleq \{\, \tau,\ \identifiable{schedule},\ \identifiable{cancel},\ \identifiable{suspend},\ \identifiable{resume},\ \identifiable{wait},\ \identifiable{provide},\ \identifiable{fail} \,\}
  \]
  where $\tau$ denotes an internal (silent) transition.
\end{definition}

\subsubsection{Inference Rules}

We present the transition relation via inference rules.
A transition $s \xrightarrow{l} s'$ is derivable when the premise is satisfied.

\medskip\noindent
\textbf{Scheduling and cancellation from Pending:}

\begin{mathpar}
  \inferrule{ }{\identifiable{Pending} \xrightarrow{\identifiable{schedule}} \identifiable{Scheduled}}
  \and
  \inferrule{ }{\identifiable{Pending} \xrightarrow{\identifiable{cancel}} \identifiable{Cancelled}}
\end{mathpar}

\medskip\noindent
\textbf{Scheduling and cancellation from Scheduled:}

\begin{mathpar}
  \inferrule{ }{\identifiable{Scheduled} \xrightarrow{\tau} \identifiable{Running}}
  \and
  \inferrule{ }{\identifiable{Scheduled} \xrightarrow{\identifiable{cancel}} \identifiable{Cancelled}}
\end{mathpar}

\medskip\noindent
\textbf{Transitions from Running:}

\begin{mathpar}
  \inferrule{ }{\identifiable{Running} \xrightarrow{\identifiable{wait}} \identifiable{Waiting\_input}}
  \and
  \inferrule{ }{\identifiable{Running} \xrightarrow{\identifiable{suspend}} \identifiable{Suspended}}
  \and
  \inferrule{ }{\identifiable{Running} \xrightarrow{\identifiable{cancel}} \identifiable{Cancelled}}
  \and
  \inferrule{ }{\identifiable{Running} \xrightarrow{\identifiable{fail}} \identifiable{Failed}}
  \and
  \inferrule{ }{\identifiable{Running} \xrightarrow{\tau} \identifiable{Completed}}
\end{mathpar}

\medskip\noindent
\textbf{Transitions from Waiting\_input:}

\begin{mathpar}
  \inferrule{ }{\identifiable{Waiting\_input} \xrightarrow{\identifiable{provide}} \identifiable{Running}}
  \and
  \inferrule{ }{\identifiable{Waiting\_input} \xrightarrow{\identifiable{cancel}} \identifiable{Cancelled}}
  \and
  \inferrule{ }{\identifiable{Waiting\_input} \xrightarrow{\identifiable{fail}} \identifiable{Failed}}
  \and
  \inferrule{ }{\identifiable{Waiting\_input} \xrightarrow{\tau} \identifiable{Completed}}
\end{mathpar}

\medskip\noindent
\textbf{Transitions from Suspended:}

\begin{mathpar}
  \inferrule{ }{\identifiable{Suspended} \xrightarrow{\identifiable{resume}} \identifiable{Scheduled}}
  \and
  \inferrule{ }{\identifiable{Suspended} \xrightarrow{\identifiable{resume}} \identifiable{Running}}
  \and
  \inferrule{ }{\identifiable{Suspended} \xrightarrow{\identifiable{cancel}} \identifiable{Cancelled}}
  \and
  \inferrule{ }{\identifiable{Suspended} \xrightarrow{\identifiable{fail}} \identifiable{Failed}}
  \and
  \inferrule{ }{\identifiable{Suspended} \xrightarrow{\tau} \identifiable{Completed}}
\end{mathpar}

\medskip\noindent
\textbf{No transitions from terminal states:}

\begin{mathpar}
  \inferrule{ }{\identifiable{Completed} \xrightarrow{l} s'} \quad \text{(for no } l, s')
  \and
  \inferrule{ }{\identifiable{Failed} \xrightarrow{l} s'} \quad \text{(for no } l, s')
  \and
  \inferrule{ }{\identifiable{Cancelled} \xrightarrow{l} s'} \quad \text{(for no } l, s')
\end{mathpar}

\subsubsection{Formal Properties}

We state the key properties of the LTS. No formal proofs are required.

\begin{prop}[Terminality]
  \label{prop:terminality}
  For all $s \in \mathcal{T}$ and all $l \in \mathcal{L}$, there does not exist $s'$
  such that $s \xrightarrow{l} s'$.
\end{prop}

\begin{prop}[No Self-Transitions]
  \label{prop:no-self}
  For all $s \in \mathcal{S}$ and all $l \in \mathcal{L}$, there is no transition
  $s \xrightarrow{l} s$.
\end{prop}

\begin{prop}[Deterministic Guards]
  \label{prop:det-guards}
  The validation function $\identifiable{validate} : \mathcal{S} \times \mathcal{S} \rightarrow \identifiable{result}$
  is total. For any source state $s$ and target state $s'$, $\identifiable{validate}(s, s')$
  returns either $\identifiable{Ok}$ or $\identifiable{Error}$ without raising an exception.
\end{prop}

\begin{prop}[Maximal Execution Path]
  \label{prop:max-path}
  Any execution trace from $\identifiable{Pending}$ to a terminal state in $\mathcal{T}$
  contains at most 7 transitions.
\end{prop}

\begin{prop}[Type Safety]
  \label{prop:type-safety}
  The OCaml type system enforces exhaustive pattern matching on the
  $\identifiable{task\_status}$ and $\identifiable{error\_category}$ variants.
  Consequently, no invalid transition can go unhandled at compile time.
\end{prop}

\begin{prop}[Exhaustiveness]
  \label{prop:exhaustiveness}
  The compiler flag $\identifiable{-warning}$ set to $\identifiable{@A}$ guarantees that
  any omitted case in a pattern match produces a warning (treated as an error in
  strict mode). The state machine implementation covers all possible transitions.
\end{prop}

\begin{prop}[Semantic Validity]
  \label{prop:semantic-validity}
  Even with type safety, the function $\identifiable{validate}$ must be invoked
  before any state transition to ensure semantic correctness. Type safety prevents
  \emph{syntactic} invalidity (e.g., using a state constructor that does not exist);
  $\identifiable{validate}$ prevents \emph{semantic} invalidity (e.g., attempting
  to schedule an already-completed task).
\end{prop}

\begin{prop}[Max Path Proof Sketch]
  \label{prop:max-path-proof}
  The longest possible path from $\identifiable{Pending}$ to a terminal state is:
  \[
  \identifiable{Pending} \xrightarrow{\identifiable{schedule}} \identifiable{Scheduled}
  \xrightarrow{\tau} \identifiable{Running}
  \xrightarrow{\identifiable{suspend}} \identifiable{Suspended}
  \xrightarrow{\identifiable{resume}} \identifiable{Scheduled}
  \xrightarrow{\tau} \identifiable{Running}
  \xrightarrow{\tau} \identifiable{Completed}
  \]
  This yields exactly 7 transitions. Every other valid trace is a proper prefix of
  this path or reaches a terminal state earlier.
\end{prop}

%---------------------------------------------------------------------

\subsection{Expression DSL --- Big-Step Operational Semantics}
\label{sec:expr-semantics}

We give the big-step operational semantics for the P-A-R expression DSL.
The semantics relate an expression $e$ and a context $\Gamma$ to a JSON value $v$.

\subsubsection{Syntax of Expressions}

\begin{definition}[Expression Forms]
  \label{def:expressions}
  The set of expressions $\mathcal{E}$ is defined by the following grammar:
  \[
  \begin{aligned}
    e ::=\ & \identifiable{literal}(v) \mid \identifiable{variable}(x) \mid
           \identifiable{equals}(e_1, e_2) \mid \identifiable{not\_equals}(e_1, e_2) \\
         & \mid \identifiable{greater\_than}(e_1, e_2) \mid \identifiable{less\_than}(e_1, e_2) \\
         & \mid \identifiable{and}(e_1, e_2) \mid \identifiable{or}(e_1, e_2) \mid
           \identifiable{not}(e) \\
         & \mid \identifiable{contains}(e_1, e_2) \mid \identifiable{has\_key}(e, k) \\
         & \mid \identifiable{is\_empty}(e) \mid \identifiable{matches}(e, p)
  \end{aligned}
  \]
  where $v$ ranges over JSON literals, $x$ over variable names (dot-notation paths),
  $k$ over string keys, and $p$ over regular expression patterns.
\end{definition}

The evaluation context $\Gamma : \mathcal{V} \rightarrow \mathcal{J}$ maps variable
names to JSON values, where $\mathcal{V}$ is the set of variable names and
$\mathcal{J}$ is the set of JSON values.

We write $\Gamma \vdash e \Downarrow v$ to denote that expression $e$ evaluates
to value $v$ under context $\Gamma$.

\subsubsection{Evaluation Rules}

\medskip\noindent
\textbf{Literal and Variable:}

\begin{mathpar}
  \inferrule{ }{\Gamma \vdash \identifiable{literal}(v) \Downarrow v}
  \and
  \inferrule{\Gamma(x) = v}{\Gamma \vdash \identifiable{variable}(x) \Downarrow v}
\end{mathpar}

\medskip\noindent
\textbf{Equality predicates:}

\begin{mathpar}
  \inferrule{
    \Gamma \vdash e_1 \Downarrow v_1 \quad
    \Gamma \vdash e_2 \Downarrow v_2 \quad
    v_1 =_{\identifiable{JSON}} v_2
  }{\Gamma \vdash \identifiable{equals}(e_1, e_2) \Downarrow \identifiable{true}}

  \and

  \inferrule{
    \Gamma \vdash e_1 \Downarrow v_1 \quad
    \Gamma \vdash e_2 \Downarrow v_2 \quad
    v_1 \neq_{\identifiable{JSON}} v_2
  }{\Gamma \vdash \identifiable{not\_equals}(e_1, e_2) \Downarrow \identifiable{true}}
\end{mathpar}

\medskip\noindent
\textbf{Relational predicates:}

\begin{mathpar}
  \inferrule{
    \Gamma \vdash e_1 \Downarrow v_1 \quad
    \Gamma \vdash e_2 \Downarrow v_2 \quad
    v_1 >_{\identifiable{JSON}} v_2
  }{\Gamma \vdash \identifiable{greater\_than}(e_1, e_2) \Downarrow \identifiable{true}}

  \and

  \inferrule{
    \Gamma \vdash e_1 \Downarrow v_1 \quad
    \Gamma \vdash e_2 \Downarrow v_2 \quad
    v_1 <_{\identifiable{JSON}} v_2
  }{\Gamma \vdash \identifiable{less\_than}(e_1, e_2) \Downarrow \identifiable{true}}
\end{mathpar}

\medskip\noindent
\textbf{Logical connectives:}

\begin{mathpar}
  \inferrule{
    \Gamma \vdash e_1 \Downarrow \identifiable{true} \quad
    \Gamma \vdash e_2 \Downarrow \identifiable{true}
  }{\Gamma \vdash \identifiable{and}(e_1, e_2) \Downarrow \identifiable{true}}

  \and

  \inferrule{
    \Gamma \vdash e_1 \Downarrow \identifiable{false}
  }{\Gamma \vdash \identifiable{and}(e_1, e_2) \Downarrow \identifiable{false}}

  \and

  \inferrule{
    \Gamma \vdash e_1 \Downarrow \identifiable{true}
  }{\Gamma \vdash \identifiable{or}(e_1, e_2) \Downarrow \identifiable{true}}

  \and

  \inferrule{
    \Gamma \vdash e_1 \Downarrow \identifiable{false} \quad
    \Gamma \vdash e_2 \Downarrow v_2
  }{\Gamma \vdash \identifiable{or}(e_1, e_2) \Downarrow v_2}

  \and

  \inferrule{
    \Gamma \vdash e \Downarrow \identifiable{false}
  }{\Gamma \vdash \identifiable{not}(e) \Downarrow \identifiable{true}}

  \and

  \inferrule{
    \Gamma \vdash e \Downarrow \identifiable{true}
  }{\Gamma \vdash \identifiable{not}(e) \Downarrow \identifiable{false}}
\end{mathpar}

\medskip\noindent
\textbf{Collection operations:}

\begin{mathpar}
  \inferrule{
    \Gamma \vdash e_1 \Downarrow v_1 \quad
    \Gamma \vdash e_2 \Downarrow v_2 \quad
    \identifiable{contains}(v_1, v_2) = \identifiable{true}
  }{\Gamma \vdash \identifiable{contains}(e_1, e_2) \Downarrow \identifiable{true}}

  \and

  \inferrule{
    \Gamma \vdash e \Downarrow v \quad
    \identifiable{has\_key}(v, k) = \identifiable{true}
  }{\Gamma \vdash \identifiable{has\_key}(e, k) \Downarrow \identifiable{true}}

  \and

  \inferrule{
    \Gamma \vdash e \Downarrow v \quad
    \identifiable{is\_empty}(v) = \identifiable{true}
  }{\Gamma \vdash \identifiable{is\_empty}(e) \Downarrow \identifiable{true}}
\end{mathpar}

\medskip\noindent
\textbf{Regex matching:}

\begin{mathpar}
  \inferrule{
    \Gamma \vdash e \Downarrow v \quad
    \identifiable{matches}(v, p) = \identifiable{true}
  }{\Gamma \vdash \identifiable{matches}(e, p) \Downarrow \identifiable{true}}
\end{mathpar}

The remaining cases (where the predicate does not hold) evaluate to
$\identifiable{false}$ by symmetry with the rules above.

\subsubsection{Resource Bounds}

The evaluator enforces two resource limits to ensure termination:

\begin{prop}[Resource Limits]
  \label{prop:resource-limits}
  The evaluator maintains a visit counter bounded by $\identifiable{max\_node\_visits} = 1000$.
  The recursive depth is bounded by $\identifiable{max\_depth} = 10$.
\end{prop}

\begin{prop}[Termination]
  \label{prop:termination}
  For any expression $e$ and context $\Gamma$, the evaluator terminates
  in at most 1000 node visits.
\end{prop}

\begin{prop}[Totality]
  \label{prop:totality}
  The evaluator is a total function: for all $e$ and $\Gamma$, either
  $\Gamma \vdash e \Downarrow v$ for some $v$, or the evaluator raises
  $\identifiable{Resource\_limit}$ (which is caught and converted to $\identifiable{Error}$).
\end{prop}

The totality theorem follows from the fact that the grammar of expressions is
context-free and the evaluator performs structural recursion over the expression
tree, decrementing the visit counter at each step. When the counter reaches zero,
$\identifiable{Resource\_limit}$ is raised.

%---------------------------------------------------------------------

\subsection{Middleware Composition --- Russian Doll Model}
\label{sec:middleware}

We formalize the middleware composition in P-A-R using the Russian Doll model,
where each middleware wraps the next, forming a nested chain of interceptors.

\subsubsection{Middleware Hook Structure}

\begin{definition}[Middleware Hook]
  \label{def:hook}
  A middleware hook $h$ is a record containing optional interceptor functions:
  \[
  h \triangleq \left\{
    \begin{aligned}
      &\identifiable{on\_before\_llm} : \identifiable{conv} \rightarrow \identifiable{conv option}, \\
      &\identifiable{on\_after\_llm}  : \identifiable{resp} \rightarrow \identifiable{resp option}, \\
      &\identifiable{on\_before\_tool} : \identifiable{call} \rightarrow \identifiable{call option}, \\
      &\identifiable{on\_after\_tool}  : \identifiable{call} \times \identifiable{result} \rightarrow \identifiable{result option}, \\
      &\identifiable{on\_error}       : \identifiable{error} \rightarrow \identifiable{result option}
    \end{aligned}
  \right\}
  \]
  where each field is $\identifiable{None}$ or $\identifiable{Some}\ f$ for a
  total function $f$.
\end{definition}

The special hook $\identifiable{id}$ where all fields are $\identifiable{None}$
is the \emph{pass-through} hook that leaves data unchanged.

\subsubsection{Composition Operator}

\begin{definition}[Composition Operator $\oplus$]
  \label{def:composition}
  For two hooks $h_1$ and $h_2$, their composition $h_1 \oplus h_2$ is defined
  pointwise on each field. For a field $\identifiable{on\_before\_llm}$:
  \[
  (h_1 \oplus h_2).\identifiable{on\_before\_llm}(c) =
  \begin{cases}
    h_1.\identifiable{on\_before\_llm}(c') & \text{if } h_2.\identifiable{on\_before\_llm}(c) = \identifiable{Some}\ c' \\
    h_1.\identifiable{on\_before\_llm}(c) & \text{if } h_2.\identifiable{on\_before\_llm}(c) = \identifiable{None} \\
    c & \text{if both are } \identifiable{None}
  \end{cases}
  \]
  The same pattern applies to all other fields. In terms of the implementation
  using $\identifiable{fold\_right}$:
  \[
  (h_1 \oplus h_2)(x) = h_1(h_2(x))
  \]
  for each interceptor, treating $\identifiable{None}$ as the identity.
\end{definition}

In the OCaml implementation, the application order follows $\identifiable{fold\_right}$:

\begin{lstlisting}[style=ocaml]
let apply_before_llm hooks conv next =
  fold_right (fun hook acc c ->
    match hook.on_before_llm with
    | Some f -> (match f c with Some c' -> acc c' | None -> acc c)
    | None -> acc c
  ) hooks next conv
\end{lstlisting}

This means the last hook in the list is applied first (innermost), and the
first hook is applied last (outermost).

\subsubsection{Monoid Properties}

We establish that middleware hooks under $\oplus$ form a monoid.

\begin{prop}[Closure]
  \label{prop:closure}
  The composition $h_1 \oplus h_2$ of any two middleware hooks $h_1, h_2$ is
  itself a middleware hook.
\end{prop}

\begin{prop}[Associativity]
  \label{prop:associativity}
  For all hooks $h_1, h_2, h_3$:
  \[
  (h_1 \oplus h_2) \oplus h_3 = h_1 \oplus (h_2 \oplus h_3)
  \]
\end{prop}

\begin{prop}[Proof sketch for Associativity]
  Function composition is associative: for any $x$,
  \[
  ((h_1 \oplus h_2) \oplus h_3)(x) = (h_1 \oplus h_2)(h_3(x))
  = h_1(h_2(h_3(x))) = h_1 \oplus (h_2 \oplus h_3)(x)
  \]
  The $\identifiable{fold\_right}$ implementation preserves this order:
  applying hooks in reverse list order for ``before'' phases and forward order
  for ``after'' phases yields the same composed function regardless of
  parenthesization.
\end{prop}

\begin{prop}[Identity Element]
  \label{prop:identity}
  The pass-through hook $\identifiable{id}$ (all fields $\identifiable{None}$)
  is the identity element: for all hooks $h$,
  \[
  h \oplus \identifiable{id} = \identifiable{id} \oplus h = h
  \]
\end{prop}

\begin{prop}[Proof sketch for Identity]
  When $\identifiable{on\_before\_llm} = \identifiable{None}$ for the identity hook,
  the $\identifiable{fold\_right}$ body returns its accumulator unchanged.
  Thus composing with $\identifiable{id}$ leaves the other hook's behavior
  unmodified.
\end{prop}

\begin{prop}[Interceptor Ordering]
  \label{prop:ordering}
  Hooks execute in the following order:
  \begin{itemize}
    \item For ``before'' phases ($\identifiable{on\_before\_llm}$, $\identifiable{on\_before\_tool}$):
          $\identifiable{fold\_right}$ applies hooks in reverse list order
          (last hook is innermost).
    \item For ``after'' phases ($\identifiable{on\_after\_llm}$, $\identifiable{on\_after\_tool}$):
          hooks execute in forward list order (first hook is outermost).
    \item For $\identifiable{on\_error}$: execution follows the same forward order
          as ``after'' phases.
  \end{itemize}
\end{prop}

This ordering reflects the Russian Doll (or onion-ring) model: outer wrappers
execute their ``before'' logic first and their ``after'' logic last.

\end{document}
